\section{Limitations}\label{sec:limitations}
This study has several limitations. All circuits are run on a classical simulator, so the results do not show how the models would perform on a noisy quantum processor. To keep the simulation fast, the denoisers are trained on part of the data, and the quantum fine-tuning uses at most 8{,}000 patches. The image-level metrics come from small test batches (32 images for the main comparison and 64 for the V2 benchmark), and the FID uses 1{,}000 images. Each model was trained once at a single noise level on one dataset, so the small differences between the 4-qubit models are not statistically established. The classical layers have far more parameters than the circuit, and we did not include a parameter-matched classical replacement, so the gain cannot be credited to the circuit alone. EPQD-6Q-V2 also changes several design choices at once, so its result cannot be attributed to the qubit count. Finally, the autoencoder limits the final image quality, since no denoiser can do better than the decoded clean latent map. Repeating the experiments with several random seeds, noise levels and datasets would give stronger evidence.
